\section{The next comparison}

The exact soft factorization and bounded coarse corrections reduce
the required logarithm to the fine soft scalar/vector pair. Its
short-current-time interval \(0<s<1\) and its physical infrared
tail \(s>K^2\) both have uniformly bounded affine contributions.
Writing
\[
 h^{\rm soft}_{M,K}(s)=[t^2]\left\{
   \operatorname{Tr}_{0,E}e^{-sV(t)}
       -2\operatorname{Tr}_{0,S}e^{-sT_s(t)}\right\},
\]
we therefore have the actual reduction
\begin{equation}
 \boxed{\displaystyle
 \beta(\mathcal D_{M,K})
       =\int_1^{K^2}h^{\rm soft}_{M,K}(s)\,\frac{ds}{s}+O(1),}
 \label{soft-final-band-reduction}
\end{equation}
uniformly in preceding depth and endpoint on the stated structural
systems. For \(K=1\) the intermediate interval is empty. The
growing hard-projection vertices have been removed from this
calculation by exact algebra. The actual effective-action Hessian
remains present, with its now controlled relative physical vertices.

Two distinct ingredients remain in identifying the logarithm.
First, the soft infrared penalties must be compared to a specified
massive reference without discarding their gauge cancellations.
Second, the actual effective-action vertices must be compared to
the reference vertices with constants uniform in preceding depth.
The previously proved nonlinear Ward extraction matches their
marginal momentum jets through field degree four; a summable
remainder estimate for the determinant is still required.
A sufficient concrete target, with \(\delta>0\), is
\[
 \left|h^{\rm soft}_{M,K}(s)-11\alpha_g\right|
       \le C\left(s^{-\delta}+(s/K^2)^\delta\right),
       \qquad 1\le s\le K^2.
\]
Its logarithmic-time integral gives the required
\(22\alpha_g\log K+O(1)\). This displayed comparison is a target,
not an estimate proved in this notebook.

For a fixed logarithmically divergent one-loop lattice integral,
Adams--Lee prove a universal logarithmic term with a finite remainder.
Their theorem provides a useful reference-integral endpoint once
its hypotheses and normalization are matched. It does not by itself
apply uniformly to our growing block fibers and varying effective
actions. The model-specific comparison must establish that interface.

A positive actual block drift would activate the finite transfer
and tuning results in Notebook 32. The full existence argument also
requires the recorded measure matching, mixed and large-field
histories, sourced continuum limits, nontriviality, Euclidean
covariance, positivity and reconstruction. These remain part of the
original goal. The spectral gaps here are Gaussian fluctuation
estimates, not a Yang--Mills mass-gap theorem.

\begin{thebibliography}{9}
\bibitem{priorsoft}
\emph{Toward a four-dimensional Yang--Mills construction},
Notebooks 32--35 and the earlier classical support and Ward
extraction modules in this workspace.
\bibitem{adamslee}
D.~H.~Adams and W.~Lee,
\emph{Structure of logarithmically divergent one-loop lattice
Feynman integrals}, Phys.\ Rev.\ D \textbf{77} (2008), 045010,
\href{https://arxiv.org/abs/0709.0781}{arXiv:0709.0781}.
\bibitem{sourcesoft}
E.~Sussman and collaborators,
\href{https://github.com/ethanwsussman/Yang--Mills-LLM-experiment}
{Yang--Mills-LLM-experiment}, deposited accounts pinned at
commit \texttt{e392ef45013a581664bab76b90867170b30fe10e}.
\end{thebibliography}
\end{document}
